\section{BH1: The settlement price}\label{sec:bh1}\label{sec:distortion}

The first hypothesis concerns the price the window produces. Its primary measure is the
distance between the settlement VWAP and the mid at the start of the window, which compares the
settlement price with the price at which the window opened.

\subsection{The window on the order book sample}

On the order book sample the settlement VWAP lies \AbsDriftBpsExpiry{} basis points from the
mid at the start of the window on expiry sessions and \AbsDriftBpsControl{} on control sessions. The difference is
not significant ($p = \SOnePValue{}$), and the primary test of BH1 is not rejected. The
simplest form of the same comparison, the move from the window's first minute to its last, is
also not significantly larger ($p = \MovPValue{}$).

Three supporting measures do differ. The settlement VWAP ends further from the last traded
price, \AbsReversalBpsExpiry{} basis points against \AbsReversalBpsControl{}
($p = \STwoPValue{}$); the final minute's VWAP diverges further from the window's
($p = \SThreePValue{}$); and less volume is traded within a basis point of the running VWAP
($p = \SFourPValue{}$). The first of these is larger in securities with derivatives than in
the placebo group, by \DidAbsreversalbpsPlacebo{} basis points
($p = \DidAbsreversalbpsPlaceboP{}$). All three describe a more volatile window rather than
trading aimed at the average. Volume does not concentrate in the final minute
($p = \SFivePValue{}$) and is, if anything, spread more evenly across the window: the
concentration of volume across its minutes falls, against the stated direction
($d = \HTwoNineEffectSizeCohenD{}$). A participant moving an average would find it cheapest to
trade at its end, and the volume does not go there.

\subsection{The window across the year}


On the full-year sample, Table~\ref{tab:activity} reports how the settlement window on each type of session differs
from the same security's window on ordinary sessions. On an ordinary session the window
accounts for about \ActWindowshareLiquidNormal{}\% of the session's volume in the liquid
group, and similar shares in the other two.

\begin{table}[htbp]
\centering
\small
\caption{Activity in the settlement window by type of session. Each entry is the coefficient
on an indicator for the session type in a regression with security fixed effects, estimated
separately for each group, and measures the difference from ordinary sessions of the same
security. The last column is the difference between the securities with derivatives and the
placebo group, estimated with security and session fixed effects. Standard errors clustered
by session in parentheses. $^{*}$, $^{**}$ and $^{***}$ denote significance at the ten, five
and one per cent levels.}
\label{tab:activity}
\resizebox{\textwidth}{!}{\input{generated/tables/activity}}
\end{table}

On monthly expiries, trading is drawn into the window in securities with derivatives and not
in the others. The window's share of session volume rises by \ActWindowshareLiquidMonthlyexpiry{}
percentage points in the liquid group and by \ActWindowshareIlliquidMonthlyexpiry{} points
in the illiquid group, while the change in the placebo group, \ActWindowsharePlaceboMonthlyexpiry{}
points, is not distinguishable from zero. The difference between the derivatives groups and
the placebo group is \ActWindowshareDidMonthlyexpiry{} points, with a standard error of
\ActWindowshareDidMonthlyexpirySe{}. The calendar controls do not reproduce this pattern. On weekly
index expiries the window's share of volume does not change in any group. On month ends it
rises sharply in all three groups, by \ActWindowshareLiquidMonthend{} points in the liquid
group and by \ActWindowsharePlaceboMonthend{} points in the placebo group, which reflects the
closing demand of index and portfolio rebalancing; the additional rise in the securities with
derivatives, \ActWindowshareDidMonthend{} points, is smaller than on monthly expiries.

The settlement price also moves further from the pre-window price on monthly expiries. For
the securities with derivatives the absolute move increases by
\ActAbsdriftDerivativeMonthlyexpiry{} basis points from an average of
\ActAbsdriftDerivativeNormal{}, and the final five minutes of the window end
\ActAbsterminalDerivativeMonthlyexpiry{} basis points further from the settlement price.
Neither increase is specific to the derivatives groups. The placebo group shows increases of
similar size, and the differences between the two, \ActAbsdriftDidMonthlyexpiry{} and
\ActAbsterminalDidMonthlyexpiry{} basis points, are small and insignificant. Figure~\ref{fig:daytypes}
shows the same comparison in levels.

\begin{figure}[htbp]
  \centering
  \includegraphics[width=\textwidth]{day_types.pdf}
  \caption{Share of session volume traded in the settlement window, and the absolute move
    from the pre-window price to the settlement price, by group and type of session. Means
    over security-sessions.}
  \label{fig:daytypes}
\end{figure}


\subsection{Reversal of the stock settlement price}

Table~\ref{tab:reversal} estimates equation~\eqref{eq:reversal}. On ordinary sessions the move
over the window is not reversed the next morning: the coefficient for the securities with
derivatives is $\RevDerivativeDrift{}$, with a standard error of \RevDerivativeDriftSe{}. On
monthly expiries a part of the move is reversed. The additional coefficient is
$\RevDerivativeDriftxmonthlyexpiry{}$ (standard error \RevDerivativeDriftxmonthlyexpirySe{}),
so that about a quarter of the window's move is undone by the next morning, and the
randomization $p$-value is \RevDerivativeDriftxmonthlyexpiryRip{}. The corresponding
coefficients for weekly index expiries and month ends are smaller and not significant.

\begin{table}[htbp]
\centering
\small
\caption{Reversal of the settlement price. Regressions of the return from the settlement price
to the volume-weighted average price of the next session's first thirty minutes on the move
over the window and its interactions with the session types, with security and session fixed
effects. The last column is the difference between the securities with derivatives and the
placebo group. Standard errors clustered by session in parentheses. The randomization
$p$-value compares the monthly-expiry coefficient with its distribution over \RiDraws{}
samples in which the expiries are replaced by one ordinary session from each month.}
\label{tab:reversal}
\resizebox{\textwidth}{!}{\input{generated/tables/reversal}}
\end{table}

The reversal is not, however, confined to securities with an expiring contract. In the placebo
group the additional coefficient on monthly expiries is $\RevPlaceboDriftxmonthlyexpiry{}$
(randomization $p$-value \RevPlaceboDriftxmonthlyexpiryRip{}), and the difference between the
two sets of securities, $\RevDidDriftxmonthlyexpiry{}$, is close to zero (randomization
$p$-value \RevDidDriftxmonthlyexpiryRip{}). A settlement price displaced by trading aimed at
the settlement should be reversed where there is a settlement to influence and not elsewhere.
The pattern found here is a property of the monthly expiry session as a whole, shared by
securities on which nothing settles.

Figure~\ref{fig:path} shows the same result as a path. For each security and session the
prices through the window and at the next morning's open are measured relative to the average
of all securities on that session, and oriented in the direction of the window's move. On
ordinary sessions the move is kept overnight; on monthly expiries part of it is given back, to
a similar extent in all three groups.

\begin{figure}[htbp]
  \centering
  \includegraphics[width=\textwidth]{window_path.pdf}
  \caption{Average path of the price through the settlement window and at the next morning's
    open, relative to the pre-window price and to the average of all securities on the same
    session, oriented so that the window's move is positive.}
  \label{fig:path}
\end{figure}

\subsection{Is the next morning's reversal news?}\label{sec:news}

Information that arrives after the close moves the next morning's price, and could produce the
fall in Figure~\ref{fig:path} without any role for the settlement. Information that is not
anticipated is unrelated to the direction of the window's move, so on average it should not
appear as a reversal of it. In a sample of twelve expiries, however, a few days with large
news can create such a relation by chance, and market-wide news moves stocks with a high
exposure to the market more, which a session fixed effect does not remove. Four checks
address these possibilities. Stock-days with a sign of news the next morning are removed:
those on which the next session's volume is more than twice the security's median, or on
which the opening gap is among the largest twentieth of the security's own gaps. The next
morning's return is adjusted for market news by subtracting the security's beta, estimated
from close-to-close returns over the year, times the average next-morning return of the
universe. The expiry of February, on the day Russia's invasion of Ukraine began, when the
market fell sharply during the session and recovered the next morning, is dropped. And the
reversal is estimated for each expiry separately.

\begin{table}[htbp]
\centering
\small
\caption{Additional reversal of the window's move on monthly expiries under the checks for
news. Coefficient on the window's move interacted with the monthly expiry in the regression of
Table~\ref{tab:reversal}, with security and session fixed effects; standard errors clustered by
session in parentheses, $p$-values in brackets.}
\label{tab:news}
\begin{tabular}{lcc}
\toprule
 & Securities with derivatives & Placebo \\
\midrule
Baseline & $\NewsDerivativeBase{}$ (\NewsDerivativeBaseSe{}) [\NewsDerivativeBaseP{}] &
  $\NewsPlaceboBase{}$ (\NewsPlaceboBaseSe{}) [\NewsPlaceboBaseP{}] \\
Without stock-days with signs of news & $\NewsDerivativeNonews{}$ (\NewsDerivativeNonewsSe{}) [\NewsDerivativeNonewsP{}] &
  $\NewsPlaceboNonews{}$ (\NewsPlaceboNonewsSe{}) [\NewsPlaceboNonewsP{}] \\
Adjusted for market exposure & $\NewsDerivativeBeta{}$ (\NewsDerivativeBetaSe{}) [\NewsDerivativeBetaP{}] &
  $\NewsPlaceboBeta{}$ (\NewsPlaceboBetaSe{}) [\NewsPlaceboBetaP{}] \\
Without the February expiry & $\NewsDerivativeNofeb{}$ (\NewsDerivativeNofebSe{}) [\NewsDerivativeNofebP{}] &
  $\NewsPlaceboNofeb{}$ (\NewsPlaceboNofebSe{}) [\NewsPlaceboNofebP{}] \\
Adjusted, without signs of news & $\NewsDerivativeBetanonews{}$ (\NewsDerivativeBetanonewsSe{}) [\NewsDerivativeBetanonewsP{}] &
  $\NewsPlaceboBetanonews{}$ (\NewsPlaceboBetanonewsSe{}) [\NewsPlaceboBetanonewsP{}] \\
\midrule
Move from the morning to the window & $\NewsDerivativeEarlier{}$ (\NewsDerivativeEarlierSe{}) [\NewsDerivativeEarlierP{}] &
  $\NewsPlaceboEarlier{}$ (\NewsPlaceboEarlierSe{}) [\NewsPlaceboEarlierP{}] \\
\bottomrule
\end{tabular}
\end{table}

Table~\ref{tab:news} reports the results. Removing the stock-days with signs of news reduces
the additional reversal in the securities with derivatives from $\NewsDerivativeBase{}$ to
$\NewsDerivativeNonews{}$, adjusting for market exposure reduces it to
$\NewsDerivativeBeta{}$, and with both it is $\NewsDerivativeBetanonews{}$ with a $p$-value of
\NewsDerivativeBetanonewsP{}. In the placebo group the same steps reduce it from
$\NewsPlaceboBase{}$ to $\NewsPlaceboBetanonews{}$ ($p = \NewsPlaceboBetanonewsP{}$).
Figure~\ref{fig:newsexpiry} shows that the reversal is carried by a few expiries:
\NewsExpiriesInrange{} of the twelve lie within the range of the middle ninety per cent of
the \NewsNormalSessions{} ordinary sessions, and the average over the twelve, compared with the
averages over one ordinary session drawn from each month, has a randomization $p$-value of
\NewsSessionriDerivative{}. The largest reversals follow the expiries of February, March,
August and September.

\begin{figure}[htbp]
  \centering
  \includegraphics[width=\textwidth]{overnight_by_expiry.pdf}
  \caption{Slope of the next morning's return on the window's move, estimated separately for
    each session. The histogram shows the ordinary sessions and the vertical lines the monthly
    expiries.}
  \label{fig:newsexpiry}
\end{figure}

\begin{figure}[htbp]
  \centering
  \includegraphics[width=0.8\textwidth]{overnight_robustness.pdf}
  \caption{Additional reversal on monthly expiries under each check for news, with 95\%
    intervals.}
  \label{fig:newsrobust}
\end{figure}

The reversal is therefore best attributed to information arriving after the session, most of
it common to the market, rather than to the release of a settlement price displaced by trading.
It is as large in securities on which nothing settles, it is concentrated on a few dates, and
it is small and insignificant once market news and signs of stock-specific news are allowed
for. One feature does not fit news alone: the move from the morning to the start of the window
is reversed far less, $\NewsDerivativeEarlier{}$, than the window's move. The trade data
identify news only through its effect on volume and prices, and a record of company
announcements would allow a direct test.

\subsection{The expiring future and the cash settlement price}

The expiring stock future trades during the window alongside the cash market whose average
sets its settlement price, and its price shows how the market values the settlement. 
Under settlement by delivery, the settlement price of an expiring stock future divides the
cash flows of a long position between the final variation margin and the payment for the
delivered shares, but does not change their sum. A long position bought at price $F$ pays
$F$ for the shares whatever the settlement price turns out to be. The expiring future should
therefore be priced on the value of the shares, that is on the current cash price, until the
end of the window, rather than on the average that determines the settlement price. If it were
settled in cash against the average, it would instead converge to the average, which by the
last minutes of the window is almost entirely determined.

Table~\ref{tab:convergence} distinguishes the two. The futures price and the cash price are
both measured relative to the settlement price, and the first is regressed on the second. By
the minute beginning two minutes before the close, \ConvExpiryMTwoEightFixed{}\% of the window's volume has already
traded, so the settlement price is nearly fixed. A future priced on the settlement price would
not respond to the cash price at that point; a future priced on the shares would move with it
one for one. The slope is \ConvExpiryMTwoEight{}, and over the final five minutes
\ConvExpiryLastFive{}, both far from zero and close to one. The median distance of the futures
price from the settlement price in that minute, \ConvExpiryMTwoEightGapsettle{} basis points, exceeds
its distance from the concurrent cash price, \ConvExpiryMTwoEightGapcash{} basis points.
Figure~\ref{fig:convergence} shows the same comparison minute by minute.

\begin{table}[htbp]
\centering
\small
\caption{Pricing of the expiring stock future in the settlement window. Regression of the
futures price on the cash price, both in basis points relative to the settlement price, over
security-minutes of the window; a slope of one indicates that the future moves with the cash
price, a slope of zero that it is priced on the settlement price. Standard errors clustered by
security in parentheses. Control sessions are the trading day before each expiry, on which the
same future has one day to run.}
\label{tab:convergence}
\input{generated/tables/convergence}
\end{table}

\begin{figure}[htbp]
  \centering
  \includegraphics[width=0.9\textwidth]{convergence.pdf}
  \caption{Median absolute distance of the expiring future's price from the concurrent cash
    price and from the settlement price, by minute of the window.}
  \label{fig:convergence}
\end{figure}

The result has a direct implication for the settlement incentive. A holder of an expiring
stock future gains nothing from moving the settlement price, because the gain on the final
margin is paid back in the price of the delivered shares. The direct interest in the
settlement price that holders of large expiring positions are often assumed to have applies to
stock futures only in accounting terms.


\subsection{Option strikes}

The incentive is real for stock options, because the settlement price determines which series
are exercised and therefore which option holders take or make delivery. If the settlement price
were drawn toward strikes on expiry days, either through the hedging of option writers or
through trading by option holders, the distance between the settlement price and the nearest
strike would be smaller on the expiry day than on other days of the month.

\begin{table}[htbp]
\centering
\small
\caption{Distance of the settlement price to the nearest option strike, in units of the local
strike interval. Other sessions are the sessions of the same expiry month; regressions include
a fixed effect for each security and month. High option activity is option volume on the
expiry and control days above the median relative to the security's cash turnover. The change
in the window is the distance at the settlement price less the distance at the pre-window
price. Standard errors clustered by session in parentheses.}
\label{tab:pinning}
\resizebox{\textwidth}{!}{\input{generated/tables/pinning}}
\end{table}

Table~\ref{tab:pinning} finds no such attraction. On other sessions of the month the average
distance is \PinDistanceControlmean{} strike intervals, close to the value of one quarter
expected when prices are unrelated to strikes. On expiry days the difference is
$\PinDistanceExpiry{}$ intervals (standard error \PinDistanceExpirySe{}), which is small and
not significant, and the share of settlement prices within a tenth of an interval of a strike
does not rise. Nor does the settlement window move the price toward the nearest strike: the
change in distance over the window is $\PinChangeExpiry{}$ on expiry days relative to other
sessions. For securities with the most option activity, where the incentive is greatest, the
estimates point in the direction opposite to attraction. Figure~\ref{fig:pinning} shows the
distribution.

\begin{figure}[htbp]
  \centering
  \includegraphics[width=0.62\textwidth]{pinning.pdf}
  \caption{Distribution of the distance of the settlement price to the nearest strike on expiry
    sessions and on the other sessions of the same month. The dashed line is the uniform
    density.}
  \label{fig:pinning}
\end{figure}


\subsection{The index, where settlement is in cash}\label{sec:index}

Settlement by delivery removes the incentive to move a stock's settlement price. It does not
remove it for the index. Nifty 50 futures and options are settled in cash against the index's closing value, which is computed from the
closing prices of the constituents, that is from their settlement-window VWAPs. The index
options expired on every Thursday of the year and were the most heavily traded derivatives on
the exchange. If the settlement window is used to move a price, the index is where it pays.

The exchange's index file records, every second, the value of the Nifty 50 and of the 
 Nifty Next 50, an index of the next fifty large companies on which no derivatives were traded. The
second index is exposed to the same market and the same calendar but to no settlement, and it
serves as the placebo at the index level. The closing value of the Nifty 50 is approximated by
its average over the window; on the dates on which the published closing value was checked,
the approximation was within a few index points of it. The comparison uses all
\IdxSessions{} sessions, of which \IdxExpiries{} were index expiries.

\begin{table}[htbp]
\centering
\small
\caption{The index settlement. Regressions of each measure on indicators for index expiries,
every Thursday on which the index options expire, and for the monthly expiry in addition,
with a control for month ends. Difference is the Nifty 50 less the Nifty Next 50, which has no
derivatives. Heteroskedasticity-robust standard errors in parentheses. The randomization
$p$-value compares the index-expiry coefficient with its distribution over \RiDraws{} samples
in which each expiry is replaced by another session of the same week.}
\label{tab:index}
\resizebox{0.95\textwidth}{!}{\input{generated/tables/index}}
\end{table}

Table~\ref{tab:index} reports the results. Neither index, taken on its own, moves in a
consistent direction over the window on expiry days; the level of each is dominated by the
movement of the market as a whole. Their difference is not. On index expiries the change of the
Nifty 50 over the window, relative to the Nifty Next 50, differs from other sessions by $\IdxSpreadDriftExpiry{}$ basis points (standard error \IdxSpreadDriftExpirySe{}, randomization $p$-value
\IdxSpreadDriftExpiryRip{}). This is about \IdxSpreadDriftPoints{} index points at the year's
average level of \IdxNiftyLevel{}. It occurs on the weekly expiries, on which only the index
options settle, as much as on the monthly ones: the additional coefficient for monthly expiries
is $\IdxSpreadDriftMonthly{}$ and insignificant. The Nifty 50 ends the window below the 
 Nifty Next 50 on \IdxNegshareExpiry{}\% of index expiries, against \IdxNegshareOther{}\% of other
sessions. The movement takes place within the window: from its first minute to its last the
difference changes by $\IdxSpreadMoveExpiry{}$ basis points relative to other sessions
(randomization $p$-value \IdxSpreadMoveExpiryRip{}). Figure~\ref{fig:indexpath} shows the path.

\begin{figure}[htbp]
  \centering
  \includegraphics[width=0.72\textwidth]{index_path.pdf}
  \caption{The Nifty 50 relative to the Nifty Next 50 through the settlement window and at the
    next morning's open, relative to the pre-window level, on index expiries and on other
    sessions. Means with 95\% intervals.}
  \label{fig:indexpath}
\end{figure}

Figure~\ref{fig:indexexpiries} shows the move on each of the year's index expiries. The effect
is not the product of a few sessions: the negative moves are spread across the year, on weekly
and monthly expiries alike. Figure~\ref{fig:randomization} compares the estimate with its
randomization distribution: the share of placebo samples, in which the expiries are replaced by
other sessions of the same week, that produce a coefficient as large in magnitude is
\IdxSpreadDriftExpiryRip{}.

\begin{figure}[htbp]
  \centering
  \includegraphics[width=\textwidth]{index_expiries.pdf}
  \caption{Move of the Nifty 50 relative to the Nifty Next 50 from the pre-window level to the
    settlement value, on each index expiry of the year. The dotted line is the mean over index
    expiries and the dashed line the mean over sessions without an expiry.}
  \label{fig:indexexpiries}
\end{figure}

\begin{figure}[htbp]
  \centering
  \includegraphics[width=0.65\textwidth]{randomization.pdf}
  \caption{Randomization distribution of the index-expiry coefficient for the Nifty 50 less the
    Nifty Next 50: the same regression estimated \RiDraws{} times with each expiry replaced by
    another session of the same week. The solid line is the estimate on the actual expiries and
    the dotted line its mirror image.}
  \label{fig:randomization}
\end{figure}

The next morning the difference recovers by \IdxSpreadNextrevExpiry{} basis points more than
after other sessions, which is of the same size as the fall and in the opposite direction. The
estimate is imprecise (randomization $p$-value \IdxSpreadNextrevExpiryRip{}), because overnight
returns vary far more than returns within the window, so the reversal is consistent with the
fall being temporary but does not establish it.

The same pattern appears in the stocks. On index expiries the liquid group, most of whose
members are constituents of the Nifty 50, changes from the first minute of the window to the last
by $\StockidxLiquidMoveOneFiveThreeZeroIndexexpiry{}$ basis points relative to the placebo group
(standard error \StockidxLiquidMoveOneFiveThreeZeroIndexexpirySe{}). The illiquid group, most of
whose members are outside the index, moves by $\StockidxIlliquidMoveOneFiveThreeZeroIndexexpiry{}$
basis points, which is not significant.

The index settlement is therefore displaced downward on expiry days, by an amount that is small
relative to the daily movement of the index but systematic across the year's expiries, specific
to the index that carries derivatives, and concentrated in the window in which its settlement
value is fixed. The data identify participant categories but not individual traders, and they
do not show the positions held in index options, so the cause cannot be established. Two
explanations are consistent with the evidence: trading intended to lower the settlement value,
and the unwinding at the settlement price of cash hedges held against index option positions.
Either way, the settlement value of the most traded derivative in the market is systematically
lowered in the window relative to an index without derivatives.

\subsection{Index strikes}

The index options have strikes at every fifty points, and their settlement is in cash. The
distance of the Nifty 50's settlement value to the nearest multiple of fifty, in units of the
interval, averages \IdxpinNiftyDistsettleNormal{} on sessions without an expiry, the value
expected when the index is unrelated to the strikes. On index expiries it differs by
$\IdxpinNiftyDistsettle{}$ (standard error \IdxpinNiftyDistsettleSe{}), and the change in
distance over the window by $\IdxpinNiftyDistchange{}$. The downward displacement of the index
settlement documented in Section~\ref{sec:index} is therefore not a movement toward the
nearest strike. It is a movement in one direction.
